\clearpage
\section{Problem Analysis}

The distributed requirements can be decomposed into three main concerns.
\begin{itemize}
    \item The server must own the authoritative match state so that clients
    never race on shared mutable data and cannot bypass rule validation.
    \item Lobby operations and game operations can happen concurrently, so
    create/join/move requests must be synchronized at the right granularity.
    \item Client notifications must be asynchronous, because remote callbacks
    can block or fail independently of the game logic.
\end{itemize}

The first requirement rules out any peer-to-peer ownership of the board. If
clients kept live mutable state, the server would no longer be the single point
of truth and move validation would become harder to reason about. The project
therefore uses a server-centric model in which each match is represented by one
remote game object.

The second requirement suggests separating concerns between matchmaking and
gameplay. The lobby is responsible for room names, room lookup, and room
lifecycle; each game instance is responsible for turn alternation, board
updates, win/draw detection, and abandonment. This separation avoids one coarse
global lock around the whole application and allows different matches to
proceed independently.

The third requirement concerns responsiveness and failure isolation. A remote
callback must never hold the critical section that protects game state, because
the client may be slow, unreachable, or already terminated. The callback path
must therefore be decoupled from state mutation, and clients must receive
immutable snapshots rather than references to live server objects.
